\documentclass[10pt,a4paper]{amsart}
\usepackage[margin=19mm]{geometry}
\usepackage{amsmath,amssymb,mathtools}
\usepackage[scr=rsfs,cal=euler]{mathalfa}
\usepackage{xfrac}
\usepackage[unicode,hidelinks,bookmarks=false]{hyperref}
\newcommand{\ie}{i.e.\ }
\newcommand{\cf}{cf.\ }
\newcommand{\resp}{respectively\ }
\newcommand{\Subsection}{\subsection}
\providecommand{\nobreakdash}{\nobreak\hbox{-}}
\setlength{\emergencystretch}{2em}
\multlinegap0pt
\usepackage{bm}
\usepackage{bbm}
\usepackage{dsfont}
\usepackage{xspace}
\usepackage{cancel}
\usepackage{mathtools}
\usepackage{amsthm}
\makeatletter
\@ifundefined{theorem}{\newtheorem{theorem}{Theorem}}{}
\@ifundefined{lemma}{\newtheorem{lemma}[theorem]{Lemma}}{}
\makeatother
\title[Nearly tight bounds for induced subdivisions]{Nearly tight bounds for induced subdivisions}
\author{Zach Hunter, Aleksa Milojevi\'c, Patryk Morawski, Benny Sudakov}
\date{Source: arXiv:2607.06444v1 (CC BY 4.0)}
\begin{document}
\maketitle

\noindent\emph{Source and licence.} Nearly tight bounds for induced subdivisions. Version arXiv:2607.06444v1 (CC BY 4.0), distributed under the Creative Commons Attribution 4.0 International licence, as recorded in the arXiv licence metadata for this exact version and shown on the abstract page https://arxiv.org/abs/2607.06444v1. Full rights notice: licenses/arxiv-2607.06444-CC-BY-4.0.txt.\par\smallskip

\noindent\textit{Editorial scope.} Counted unit: the Lemma whose source label is \texttt{lemma:collection\_of\_paths} in the article (source0.tex lines 630--639, source label \texttt{lemma:collection\_of\_paths}), delivered together with the article's own complete original proof of it (source0.tex lines 649--697, one \texttt{proof} environment of 49 lines). No mathematics is written, repaired or re-derived: statement and proof are the article's own text line for line. The proof is detached from the statement in the source: the article states the result at lines 630--639 and prints its proof at lines 649--697, and the proof environment's own header names the statement, so the association is the article's own. Required context retained uncounted: lemma:drifting\_away\_paths (lines 577--582). Every definition, display and label of those lines is retained unchanged. Omitted from this package and not redistributed: 1--576, 583--629, 640--648, 698--1058. Nothing inside a retained excerpt is omitted or altered: every retained body is delivered exactly as the article's lines are written, so no omission receipt of any kind accompanies this package. Citations: the retained excerpts carry no citation command at all, so no bibliography entry of the article is transcribed and no reference is redistributed. Reference adaptation: none was needed and none was made; every reference inside the delivered unit resolves inside it, so no unresolved reference or citation is produced. Printed slips kept literal: no printed word, symbol, hypothesis or number of the counted statement or of its proof was altered. Layout and class: the article's own class is \texttt{article} with packages including inputenc, fontenc, graphicx, amsmath, amsthm, amsfonts, xcolor, amssymb, url, hyperref, breakurl, fullpage, esvect, mathtools; the wrapper is the lane's pinned \texttt{amsart} editorial wrapper at 10pt on A4, which restates the theorem-like environments the retained text needs and adds the article's own macro definitions that the retained text needs (none), with extra wrapper packages (bm, bbm, dsfont, xspace, cancel, mathtools). The delivered numbering is the unit's own. Assets: no retained span contains a figure environment, a graphics-inclusion command, a PDF-page inclusion, a table or a TikZ picture; no image file is copied into this package, the package declares no asset dependency and no image file is redistributed with it. Licence: version arXiv:2607.06444v1 of this article is distributed under the Creative Commons Attribution 4.0 International licence, as recorded in the arXiv licence metadata for this exact version and shown on the abstract page https://arxiv.org/abs/2607.06444v1. A full rights notice is delivered at licenses/arxiv-2607.06444-CC-BY-4.0.txt. Characters: every retained excerpt is pure ASCII, so the delivered engine, XeLaTeX with its default Latin Modern text font, reports no missing character.
\begin{lemma}\label{lemma:drifting_away_paths}
    Let $r > 0$ and let $G$ be an $r$-expander.
    Let $U, U_1, U_2 \subseteq V(G)$ be sets such that $|U_1|, |U_2| \geq r$ and $|U| \leq r/10^3$.
    Let moreover $F \subseteq V(G)$ be a set of forbidden vertices such that $|B^{(\ell)}(U_1 \cup U_2) \cap F| \leq r(\ell+1)^3/10^{15}$ for each $\ell \geq 0$.
    Then, there exist paths $P_1, P_2$ in $G$ which are disjoint from $F$ and drifting-away from $U$ such that $P_1P_2^{-1}$ is a path from $U_1$ to $U_2$ (here, $P_1P_2^{-1}$ denotes the concatenation of the path $P_1$ with the reverse of $P_2$).
\end{lemma}

\begin{lemma}\label{lemma:collection_of_paths}
    Let $h, r \geq 1$, $a \geq 16h$ and let $G$ be an $r$-expander for $r \geq 10^{23}h^2\Delta(G)\log^3\Delta(G)$.
    Let moreover $U \subseteq V(G)$ and  $S_1, \dots, S_{2h} \subseteq U$ be disjoint sets such that for each $i \in [2h]$
    \begin{enumerate}
        \item $|U| \leq r/(10^{16})$ and $|S_i| \geq a$, 
        \item any subset $S \subseteq S_i$ with $|S| \geq |S_i| /2$ has $|N(S)| \geq r$, \emph{and}
        \item for each $v \in \left(V(G) \setminus U\right) \cup S_1 \cup \dots \cup S_{2h}$ we have $|N(v) \cap (S_1\cup \dots \cup S_{2h})| \leq ha/(10^{8}h^2\log^3\Delta(G))$.
    \end{enumerate}
    Then, there exists a subset $I \subseteq [2h]$ of size $h$ and a collection of disjoint paths $\{P_{i, j} : \{i, j\} \in \binom{I}{2}\}$ such that $P_{i, j}$ is a path from $S_i$ to $S_j$, $|P_{i, j} \cap U| =2$ and the only edges spanned by $\bigcup_{i, j} V(P_{i, j})$ are $\bigcup_{i, j} E(P_{i, j})$.
\end{lemma}

\begin{proof}[Proof of Lemma~\ref{lemma:collection_of_paths}.]
    We let $G$ together with sets $S_1, \dots, S_{2h} \subseteq U \subseteq V(G)$ satisfy the conditions of the lemma and write $d = d(G)$.
    We want to find the paths connecting $S_1, \dots, S_{2h}$ one by one.
    During the process, we discard some of the sets if they lose too many vertices.
    We need to argue that we can find a suitable path in each step and that we discard at most $h$ of the sets, so that at least $h$ sets remain at the end of our process.

    More formally, for each pair $\{i, j\} \in \binom{[2h]}{2}$ we want to find an induced path $P_{i,j}$ such that
    \begin{enumerate}
        \item $P_{i, j}$ is a path whose only vertices in $U$ are its endpoints $x \in S_i$ and $x' \in S_j$ and which satisfies $|B^{(\ell)}(U)\cap V(P_{i, j})|\leq 2(20\log \Delta(G))^3 (\ell+1)^3+2$ for all $\ell\geq 0$, \emph{and}
        \item for $\{i, j\} \neq \{i', j'\}$, there are no edges crossing between $P_{i', j'}$ and $P_{i, j}$.
    \end{enumerate}
    We note that these conditions guarantee that the collection of paths we find is an induced collection of disjoint paths in $G$.

    We attempt to find the paths one by one --- in each step we pick a not-yet connected pair $\{i, j\}$ and find the path $P_{i, j}$.
    We let $\mathcal{P}_q$ denote the set of $q$ paths found up until step $q$, let $V_q = \bigcup_{P \in \mathcal{P}_q} P$ and $F_q$ be the set of neighbours of some $v\in V_q$ together with $V_q$ itself. When finding the next path, we make sure that it is disjoint from $F_q$, which makes sure that the second condition is satisfied.

    We say that a set $S_i$ becomes overused at step $q$ if $|S_i \cap F_q| \geq |S_i|/2$. If a set $S_i$ becomes overused, we immediately give up on it --- we will no longer find paths connecting it to the other vertices.

    Before arguing that we can indeed find a suitable path in each step, we first show that at any point the number of overused sets is at most $h$ --- and hence at the end of the process there will be at least $h$ of them left for us to define the set $I$.

    By the first property of the paths $P_{i, j}$, for each $\ell \geq 0$ the number of vertices in $V_q$ at distance at most $\ell$ from $U$ is at most
    \[
        2q(\ell+1)^3  \cdot 8000\log^3\Delta(G)  + 4q \leq 10^5(\ell+1)^3h^2\log^3\Delta(G).
    \]
    
    Note that by the third condition in the lemma and our bound on the number of vertices in $V_q$ at distance at most $1$ to $U$, we have
    \[
        |F_q \cap \left(S_1 \cup \dots \cup S_{2h} \right)| \leq 2q + \left(2q + 8\cdot 10^{5}h^2\log^3\Delta(G)\right) \cdot \frac{ha}{10^{8}h^2\log^3\Delta(G)} \leq ha/2.
    \]
    Since $S_i \subseteq U$, $|S_i| \geq a$ for all $i \in [2h]$ and the sets $S_i$ are pairwise disjoint, we can bound the number of overused sets at any step by $|F_q \cap \left(S_1 \cup \dots \cup S_{2h} \right)|/(a/2) \leq h$.

    Suppose that for some $0\leq q < \binom{2h}{2}$ we have found a collection $\mathcal{P}_{q}$ of paths satisfying the above conditions and let $\{i, j\} \in \binom{[2h]}{2}$ be a pair such that both $S_i$ and $S_j$ are not overused at step $q$ and that we have not found the path $P_{i, j}$ yet.
    Since $S_i$ is not overused, we have $|S_i \setminus F_{q}|\geq |S_i|/2$, and so taking $U_i = N(S_i\backslash F_q)$ we get that $|U_i| \geq r$ by our second assumption in the lemma.
    We define $U_j$ similarly for $S_j$ and again get that $|U_j| \geq r$.

    We now apply Lemma~\ref{lemma:drifting_away_paths} with the sets $U_i$ and $U_j$ and $F = F_{q} \cup U$.
    We note that, using the above observations, for each $\ell \geq 0$ the number of vertices in $F$ at distance at most $\ell$ from $U_i \cup U_j$ is at most
    \[
        |U| + 10^5(\ell+2)^3h^2\Delta(G)\log^3\Delta(G) \leq \frac{r}{10^{16}} + \frac{r + r\ell^3}{10^{16}} \leq \frac{r + r\ell^3}{10^{15}},
    \]
    where we used our assumption on $r$ and the size of $U$.
    In particular, the premise of Lemma~\ref{lemma:drifting_away_paths} is satisfied.
    We thus get paths $P_1, P_2$ that are disjoint from $F$ and drifting-away from $U$ such that $P_1P_2^{-1}$ is a path from $U_1$ to $U_2$.
    We pick $x_1 \in N(P_1) \cap S_i\backslash F_q$ and $x_2 \in N(P_2) \cap S_j\backslash F_q$.
    Since the induced subgraph on $V(P_1)\cup V(P_2)\cup \{x_1, x_2\}$ is connected, we may take $P_{i, j}$ to be the shortest path between $x_1$ and $x_2$ in this subgraph, and note it must be induced due to minimality. Finally, this path satisfies the first property since $P_1$ and $P_2$ are drifting away from $U$ and so \[|B^{(\ell)}(U)\cap V(P_{i, j})|\leq |B^{(\ell)}(U)\cap V(P_1)|+|B^{(\ell)}(U)\cap V(P_2)|+2\leq 2(20\log \Delta(G))^3 (\ell+1)^3+2.\]
    
    At the end of the process, we have found a suitable path between all pairs of non-overused sets.
    Since there are at least $h$ non-overused sets, we can pick $I \subseteq [2h]$ of size $h$ such that the corresponding subset of the paths satisfies the conditions of the lemma.
\end{proof}

\end{document}
